\subsection{Perturbation of Riesz endomorphisms}

Lemma 1. --- Let E be a Banach space. Let p and q be continuous projectors on E such that \(\|q - p\| < 1\) . Then p induces an isomorphism of \(\operatorname{Im}(q)\) onto \(\operatorname{Im}(p)\) , and \(1_{E} - p\) induces an isomorphism of \(\operatorname{Ker}(q)\) onto \(\operatorname{Ker}(p)\) .

Consider the continuous linear maps \(u\colon x\mapsto p(x)\) from \(\operatorname{Im}(q)\) into \(\operatorname{Im}(p)\) and \(v\colon y\mapsto q(y)\) from \(\operatorname{Im}(p)\) into \(\operatorname{Im}(q)\) . For every \(x\in\operatorname{Im}(q)\) , we have \(x=q(x)\) , whence

\[
\begin{array}{c} (q - p) ^ {2} (x) = q ^ {2} (x) - q (p (x)) - p (q (x)) + p ^ {2} (x) \\ = x - q (p (x)) = x - v (u (x)). \end{array}
\]

It follows that we have \(\|1_{E}-v\circ u\|\leqslant\|q-p\|^{2}<1\) . By cor. 1 of I, p. 22, \(v\circ u\) is an automorphism of \(\operatorname{Im}(q)\) . One proves similarly that \(u\circ v\) is an automorphism of \(\operatorname{Im}(p)\) . This implies that u is an isomorphism from \(\operatorname{Im}(q)\) onto \(\operatorname{Im}(p)\) , whence the first assertion of the lemma. The second follows from it by replacing p and q by \(1_{E}-p\) and \(1_{E}-q\) respectively.

THEOREM 2. --- Let E be a Banach space. The set \(\mathcal{R}(\mathrm{E})\) of Riesz endomorphisms of E is open in \(\mathcal{L}(\mathrm{E})\) . The map that assigns to an element of \(\mathcal{R}(\mathrm{E})\) the dimension of its nilspace is upper semicontinuous on \(\mathcal{L}(\mathrm{E})\) .

The theorem follows from the following more precise statement:

PROPOSITION 5. --- Let E be a Banach space and \(u_{0}\) a Riesz endomorphism of E. Let N (resp. I) be the nilspace (resp. the conilspace) of \(u_{0}\) . Let d denote the dimension of N. There exists an open neighborhood U of \(u_{0}\) in \(\mathcal{L}(\mathrm{E})\) and an analytic mapping \(\pi\colon\mathrm{U}\to\mathcal{L}(\mathrm{E})\) such that

\begin{enumerate}
\def\labelenumi{(\roman{enumi})}
\item
  The endomorphism \(\pi(u_0)\) is the projector with image N and kernel I;
\item
  For every \(u \in U\) , the linear map \(\pi(u)\) is a projector of rank d that belongs to the bicommutant of u and, in particular, commutes with u. Moreover, u induces an automorphism of \(\operatorname{Ker}(\pi(u))\) ;
\item
  Every element of U is a Riesz endomorphism whose nilspace has dimension \(\leqslant d\) .
\end{enumerate}

If K = C, let \(\mathrm{Sp}(u_{0})\) denote the spectrum of \(u_{0}\) with respect to the algebra \(\mathcal{L}(\mathrm{E})\) . When \(K = R\) , let \(\mathrm{Sp}(u_{0})\) denote the complex spectrum \(\mathrm{Sp}_{\mathcal{L}(\mathrm{E}_{(\mathbf{C})})}((u_{0})_{(\mathbf{C})})\) of \(u_{0}\) (I, p. 85, no. 13). By Prop. 14 of III, p. 54, the point 0 is isolated in \(\{0\} \cup \mathrm{Sp}(u_{0})\) . Let r \textgreater{} 0 be a real number such that every element of \(\mathrm{Sp}(u_{0}) - \{0\}\) has modulus \textgreater{} r. Let V denote the open set in C consisting of complex numbers of absolute value ≠ r; let f be the holomorphic function on V defined by \(f(z) = 0\) if \(|z| > r\) and \(f(z) = 1\) if \(|z| < r\) . If K = C (resp. K = R), let U' denote the set of elements u of \(\mathcal{L}(\mathrm{E})\) whose spectrum (resp. whose complex spectrum) is contained in V. The set U' is an open neighborhood of \(u_{0}\) in \(\mathcal{L}(\mathrm{E})\) and the map \(u \mapsto f(u)\) from U' into \(\mathcal{L}(\mathrm{E})\) is holomorphic (I, p. 76, prop. 10 and I, p. 85, no. 13).

Let \(u \in \mathrm{U}'\) . The endomorphism \(f(u)\) is the spectral projector \(e_0(u)\) ; it commutes with \(u\) , and \(u\) induces, by passing to subspaces, an automorphism of \(\operatorname{Ker}(e_0(u))\) (cf. III, p. 53, no. 6).

The projector \(e_{0}(u_{0})\) has image N and kernel I (III, p. 54, prop. 14); its rank is d. By Lemma 1, the set U of elements \(u \in U'\) such that \(e_{0}(u)\) is of rank d is an open neighborhood of \(u_{0}\) . The set U and the map \(\pi: u \mapsto e_{0}(u)\) from U into \(\mathcal{L}(\mathrm{E})\) satisfy conditions (i) and (ii) of the proposition.

Let \(u \in U\) . Since u induces an automorphism of \(\operatorname{Ker}(e_{0}(u))\) , which is a closed vector subspace of finite codimension of E, the endomorphism u is a Riesz endomorphism of E (III, p. 48, prop. 8); the nilspace of u is contained in the image of \(\pi(u) = e_{0}(u)\) , and has dimension \(\leqslant d\) .

